% !TEX root = main.tex


\section{Methods} \label{sec:methods}
This section describes the methodological basis to address the prediction task outlined in \Cref{sec:prediction_task} given the recorded defect severity trajectories in \Cref{sec:data}. 
We propose the Markov jump process (or \textit{continuous-time Markov chain}) for the prediction and modeling task as defect growth is measured as jumps in the state space, from one class to another.
We describe its mathematical properties with computational remarks, if any, and lay out how it is estimated and applied for prediction. In addition, we describe the evaluation methodology with respect to error metrics and reference predictors. 
We denote the complete set of rail defects by $\mathcal{S}$. The set of possible defect classifications is defined as $\mathcal{E} = \{3, 2B, 2A, 1, 0\}$, where $\mathcal{E}$ is ordered with index set denoted $\mathcal{J}$.
Further, we define a stochastic process by $\{ X^{(s)}(t) \; \vert \; t \in \mathbb R_+\}$ for each $s \in \mathcal{S}$, and let its range be given by the defect classifications such that $ X^{(s)}(t) \in \mathcal{J}, \; \forall t \in \mathbb R_+$. Hence, the state space of $X^{(s)}(t)$ is discretely valued and finite. 
The recorded endpoint data of the defect trajectories described in \Cref{sec:defect_trajectories} are denoted as $Y^{(s)} = \{ y_1^{(s)} = X^{(s)}(t_1^{(s)}), y_2^{(s)} = X^{(s)}(t_2^{(s)}) \}$ with $y_1^{(s)}$ and $y_2^{(s)}$ being the observed states and $t_1^{(s)}$ respectively $t_2^{(s)}$ their time stamps, for defect $s \in \mathcal{S}$. We define the time passed between two observations $\tau^{(s)} = t_2^{(s)} - t_1^{(s)}, \; s \in \mathcal{S}$. 







\subsection{Markov jump process} \label{sec:mjp}
The stochastic process $X^{(s)}(t) \in \mathcal{J}, \; \forall t \in \mathbb R_+$ is assumed to obey the first-order Markov property for a defect $s \in \mathcal S$. Additionally, temporal homogeneity of the process is assumed, i.e.
\begin{align}
    P\left(X^{(s)}(u+t) = j \mid X^{(s)}(u) = i \right) = P\left( X^{(s)}(t) = j \mid X^{(s)}(0) = i \right), \quad u,t\in \mathbb R_+, s \in \mathcal S, i,j \in \mathcal{J}. \label{eq:markov_property}
\end{align}
The assumption gives the process a memory-free property, meaning that the probability of transitioning from state $i$ to $j$ over a certain time span is the same regardless of the exact position in time and which prior states that $X^{(s)}$ has visited. \citet{bisgaard2024A} argue that this model assumption is reasonable as no structural or regimes in defect growth have been indicated by Bane NOR as well as from visual diagnostics of the recorded transition times.
However, spatial homogeneity is not assumed, which means that the propagation of defect severity can vary depending on the exact location of the defect in the rail network. We will expand on the notion of spatial inhomogeneity in the following.

The instantaneous rate of transition from state $i$ to state $j$ is defined as
\begin{align*}
    %\bm{A}_{i,j} = \lambda_{\mathcal{E}(i), \mathcal{E}(j)} =
    \lambda_{i,j} = 
    \lim_{\Delta t \rightarrow 0_+} \frac{P\left( X^{(s)}(\Delta t) = j | X^{(s)}(0) = i\right)}{\Delta t}, \quad i \neq j, s \in \mathcal S,
\end{align*}
for $i,j \in \mathcal{J}$. The transition rates are collected in the \textit{infinitesimal generator}, $\bm{A} \in \mathbb R^{m \times m}$, (or alternatively denoted the \textit{transition rate matrix}), such that the $(i,j)$'th element in $\bm{A}$ is given by $\lambda_{i,j}$ where $i \neq j$. Naturally, the rates of transition are nonnegative, i.e. $\lambda_{i,j} \geq 0$ for $i \neq j$. The diagonal elements are $\bm{A}_{i,i} = -\sum_{j \in \mathcal{J} | i \neq j } \lambda_{i,j}$ and reflect the total rate of leaving state $i$ \citep{bladt_nielsen}. 
To introduce exogenous dependencies to the Markov jump process, we regress the transition rate matrix on the covariates, $\bm{z}^{(s)}$, by
\begin{equation}
\begin{aligned}
    \bm{A}^{(s)} = \bm{A} \cdot g\left( \bm{\beta}^\top \bm{z}^{(s)} \right)
    %\lambda_{i,j} ( \bm{z}^{(s)} ) = \lambda_{i,j}^{(0)} \; g\left( \bm{\beta}^\top \bm{z}^{(s)} \right), \quad i,j \in \mathcal{J}
\end{aligned} \label{eq:trans_rates_covariates}
\end{equation}
with $\bm{A}$ denoting some base transition rate matrix. The function $g \, : \, \mathbb{R} \rightarrow \mathbb{R_+}$ takes any covariate dot product, $\bm{\beta}^\top \bm{z}^{(s)}$, as input and maps it to the positive real line, ensuring that the transition rates are nonnegative. We examine different choices for $g(\cdot)$, namely 
\begin{align}
    g(x) = \exp(x), \quad g(x) = \log(1 + \exp(x) ), \quad \text{and} \quad g(x) = x^2.
\end{align}
$\bm{z}^{(s)}$ contains the exogenous information described in \Cref{sec:data}. With this exogeneous information, we deem the Markov chain spatially inhomogeneous as the rail characteristics vary throughout the network as described in \Cref{sec:data}. This dependency setting is indicated by writing $\bm{A}^{(s)}$ with explicit $s$-dependence.

Let $\bm{\pi}^{(s)}(t) = \begin{bmatrix} \pi^{(s)}_1(t), \dots, \pi^{(s)}_m(t) \end{bmatrix}$ denote the time-dependent state distribution of size $1 \times m$ with $\sum_{i = 1}^m \bm{\pi}^{(s)}_i(t) = 1, \; \forall t\in \mathbb{R}, \forall s \in \mathcal{S}$. Alternatively, $\bm{\pi}^{(s)}(t)$ is denoted the \textit{transient distribution} of the Markov chain at time $t$.
The Kolmogorov forward equation, $\Dot{\bm{\pi}}^{(s)}(t) = \bm{\pi}^{(s)}(t) \bm{A}^{(s)}$, describes how the state probabilities evolve over time, given some initial probabilities \citep{bladt_nielsen}. The transient distribution can be obtained by
\begin{equation}
\begin{aligned}
                  \bm{\pi}^{(s)}(t_0 + \Delta t) = 
                  \bm{\pi}^{(s)}(t_0) \exp{ ( \Delta t \bm{A}^{(s)} ) }. \label{eq:markov_eq_solution}
\end{aligned}     
\end{equation}
\Cref{eq:markov_eq_solution} demonstrates that the transient distribution of a continuous time Markov chain is calculated as a vector-matrix product between some initial distribution of the chain, $\bm{\pi}^{(s)}(t_0)$, and the matrix exponential of the generator scaled by passed time, $\exp{ ( \Delta t \bm{A}^{(s)} ) }$. Thus, $\exp{ ( \Delta t \bm{A}^{(s)} ) }$ is a transition probability matrix, that is an $m \times m$ matrix whose $(i,j)$-indices represent the probabilities of transitioning from state $i$ to state $j$ over the time span of $\Delta t$.


\subsubsection{Temporal warping} \label{sec:warping}
Previous work on Markov jump processes applied to crack growth typically parameterizes progression in terms of calendar time, assuming that all time contributes equally to defect development \citep{bisgaard2024A}. However, for severity classes, calendar time may not contribute equally, meaning that a defect might be subject to state-dependent degradation dynamics. To model this dependency, \textit{temporal warping} is proposed.
We define
\begin{align}
    \mu_i^{(s)}(\tau ) = \int_0^\tau \pi^{(s)}_i(t) \; dt, \label{eq:cumulative_sojourn_time}
\end{align}
which is the estimated expectation of the cumulative sojourn time in state $i$ during the interval $\left[ 0,\tau \right]$ for a defect $s \in \mathcal{S}$. 
In practice, this is calculated by 
$
\bm{\mu}^{(s)}(\tau) \approx \sum_{k = 1}^{K} \Delta t \cdot \bm{\pi}^{(s)}(t_k)
$
for some large enough $K$ with $\Delta t = \tau /K$.
The warped time $\widetilde{\tau}^{(s)}$ for crack $s \in \mathcal S$, is then obtained as a weighted sum of the expected sojourn times
\begin{equation}
    \begin{aligned}
     \widetilde\tau^{(s)}(\tau) &= \sum_{i = 1}^m \xi_i  \cdot \mu_i^{(s)}(\tau ) \\
    &= \int_0^\tau \sum_{i = 1}^m \xi_i \cdot \pi^{(s)}_i \, dt \\
    &\approx \sum_{k = 1}^K \bm{\xi} \cdot \bm{\pi}^{(s)}(t_k) \cdot \Delta t,   
    \end{aligned} \label{eq:warped_time}
\end{equation}
where $\bm{\xi} \in \mathbb{R}_+^m$ is a vector of nonnegative weights. An Euler scheme calculates the transient distribution, i.e. $\bm{\pi}^{(s)}(t_k) \approx \bm{\pi}^{(s)}(t_{k-1}) + \Delta t \cdot \bm{\pi}^{(s)}(t_{k-1}) \bm{A}^{(s)}$.
If all $\xi_i = 1$, the model reduces to one without warping, i.e., $\widetilde\tau^{(s)} = \tau^{(s)}$. In this study, we set the weight for the absorbing state to $1$ to ensure that arrival at the absorbing state concludes the degradation process in warped time domain as well.
We compare models with and without temporal warping in terms of predictive accuracy. %For notational simplicity, we write $\tau'$ in both cases of standard covariate-dependent MJP and its temporally warped version.

\subsubsection{Parameterization of the infinitesimal generator} \label{sec:generator_params}
The parameterization of the infinitesimal generator (or just the \textit{generator}) determines how $X^{(s)}$ transitions between states in the Markov chain. 
%As described in \Cref{sec:data}, we model the progression of rail defects. 
Given an observation $y_1^{(s)} = X^{(s)}(t_1^{(s)})$ for defect $s \in \mathcal{S}$, a later observation $y_2^{(s)} = X^{(s)}(t_2^{(s)})$ satisfies $y_2^{(s)} \geq y_1^{(s)}$ for $t_2^{(s)} > t_1^{(s)}$, reflecting non-reversible degradation. However, between observations, the defect’s state is unobserved: $X^{(s)}(t)$ is unknown for $t \in (t_1^{(s)}, t_2^{(s)})$.

For example, if $y_1^{(s)}$ is class $2B$ and $y_2^{(s)}$ is class 1, it is not observed whether, or when, the defect entered class $2A$. By parameterizing the generator, we can either force the process to pass through $2A$ before reaching $1$, or allow direct transitions between non-neighboring states. 

Since rails can degrade both gradually and through sudden shocks, the generator’s structure influences which physical degradation mechanisms the model can capture. In the following, we introduce five different parameterizations of $\bm{A}^{(s)}$, each applied later in the prediction study to assess their impact on forecast accuracy. \Cref{tab:generators_overview} provides an overview of the generators considered in the following. Common to all is that the lower triangle of $\bm{A}^{(s)}$ is zero, reflecting that degradation is irreversible without maintenance.

\begin{table}[htb!]
\centering
\caption{Overview of generator parameterizations for crack $s \in \mathcal{S}$ where $m$ is the size of the state space and $p$ the number of covariates.}
\label{tab:generators_overview}
\begin{tabular}{@{}clcc@{}}
\toprule
Generator & Allowed jumps & Free base rates & \makecell{Total parameters} \\
\midrule
$\bm{A}_{\RN{1}}^{(s)}$ & Nearest-class & $m - 1$ & $2m - 2 + p$ \\
$\bm{A}_{\RN{2}}^{(s)}$ & Any forward jump (harmonically constrained) & $m - 1$ & $2m - 2 + p$ \\
$\bm{A}_{\RN{3}}^{(s)}$ & Up to 2 classes ahead & $2m - 3$ & $3m - 4 + p$ \\
$\bm{A}_{\RN{4}}^{(s)}$ & Up to 3 classes ahead & $3m - 6$ & $4m - 7 + p$ \\
$\bm{A}_{\RN{5}}^{(s)}$ & Any forward jump (free) & $\frac{m(m-1)}{2}$ & $\frac{m(m+1)}{2} - 1 + p$ \\
\bottomrule
\end{tabular}
\end{table}



\myparagraph{Generalized Erlang}
The first proposed parameterization assumes that a defect transition can only happen to a neighboring defect class. The generator can then be expressed as
\begin{align}
    \bm{A}_{\RN{1}}^{(s)} = \begin{bmatrix} 
    \Sigma_1 & \lambda_{3,2B} & 0 & 0 & 0 \\
    0 & \Sigma_2 & \lambda_{2B, 2A} & 0 & 0 \\
    0 & 0 & \Sigma_3 & \lambda_{2A, 1} & 0 \\
    0 & 0 & 0 & \Sigma_4 & \lambda_{1, 0} \\
    0 & 0 & 0 & 0 & 0 \\
    \end{bmatrix}, \label{eq:generator_erlang}
\end{align}
where $\Sigma_i = -\sum_{j=1, i \neq j}^m A^{(s)}_{\RN{1},i,j}$ for $i = 1, \dots, m-1$. 
\Cref{eq:generator_erlang} results in the \textit{generalized Erlang distribution} or a \textit{hypoexponential distribution} for transition times between states. This is a distribution of the sum of independent exponential random variables \citep{bladt_nielsen}. For the previously outlined observation of a defect jump from class $2B$ to class $1$ between time $t_1^{(s)}$ and $t^{(s)}_2$, \cref{eq:generator_erlang} assumes that $X^{(s)}(t^{(s)})$ has visited class $2A$ where $t^{(s)}_1 < t^{(s)} < t^{(s)}_2$.
%For brevity of notation, we have dropped the explicit $\bm{z}^{(s)}$-dependence on the $\lambda_{i,j}$'s but we note that we refer to the variables expressed in \cref{eq:trans_rates_covariates}. 
Given the parameterization in \cref{eq:generator_erlang}, the model setting has $m-1$ parameters, while it with the regression setting in \cref{eq:trans_rates_covariates} and temporal warping in \cref{sec:warping} has $2m-2 + p$ parameters, where $m$ denotes the size of the state space and $p$ denotes the number of covariates.


\myparagraph{Relaxed Erlang (harmonic-skip upper)} 
We introduce an alternative parameterization that retains the same number of free parameters as the generalized Erlang, but permits larger jumps in the state space. In this setting, longer jumps occur with lower mean transition rates, reflecting accumulated degradation over multiple intermediate classes. 
Specifically, we assume that the mean degradation time between two states is the sum of the mean degradation times through each intermediate class. For instance, we have that
\begin{align*}
    \frac{1}{\lambda_{2B, 1}} = \frac{1}{\lambda_{2B, 2A}} + \frac{1}{\lambda_{2A, 1}}
    %= \frac{\lambda_{2A, 1} + \lambda_{1, 0}}{\lambda_{2A, 1} \lambda_{1, 0}} 
    \iff \lambda_{2B,1} = \frac{\lambda_{2B, 2A} \lambda_{2A, 1}}{\lambda_{2B, 2A} + \lambda_{2A, 1}}. 
\end{align*}
Formally, this is expressed as
\begin{align}
    \frac{1}{\lambda_{i,j}} = \sum_{k = i}^{j-1} \frac{1}{\lambda_{k,k+1}}, \quad i < j. \label{eq:trans_rate_param}
\end{align}
We let $\bm{A}_{\RN{2}}^{(s)}$ denote the generator subject to the reparameterized upper triangle of the generalized Erlang in \cref{eq:generator_erlang}.
Upon determination of the upper triangle using \cref{eq:trans_rate_param}, the diagonal elements are updated accordingly so that they still satisfy $A_{\RN{2},i,i}^{(s)} = -\sum_{j=1, i \neq j}^m A_{\RN{2},i,j}^{(s)}$ for $i = 1, \dots, m$. For completeness, we provide its somewhat cumbersome algebraic expression, i.e.
\begin{align*}
  \bm{A}_{\RN{2}}^{(s)}=\begin{bmatrix}
\Sigma_1 %-\lambda_{3,2B}-\frac{\lambda_{3,2B} \lambda_{2B,2A}}{\lambda_{3,2B}+\lambda_{2B,2A}}-\frac{\lambda_{3,2B} \lambda_{2B,2A} \lambda_{2A,1}}{\left(\lambda_{2B,2A}+\lambda_{2A,1}\right) \lambda_{3,2B}+\lambda_{2B,2A} \lambda_{2A,1}}-\\\frac{\lambda_{3,2B} \lambda_{2B,2A} \lambda_{2A,1} \lambda_{1,0}}{\left(\left(\lambda_{2A,1}+\lambda_{1,0}\right) \lambda_{2B,2A}+\lambda_{2A,1} \lambda_{1,0}\right) \lambda_{3,2B}+\lambda_{2B,2A} \lambda_{2A,1} \lambda_{1,0}} 
 & \lambda_{3,2B} & \frac{\lambda_{3,2B} \lambda_{2B,2A}}{\lambda_{3,2B}+\lambda_{2B,2A}} & \frac{\lambda_{3,2B} \lambda_{2B,2A} \lambda_{2A,1}}{\left(\lambda_{2B,2A}+\lambda_{2A,1}\right) \lambda_{3,2B}+\lambda_{2B,2A} \lambda_{2A,1}} & \frac{\lambda_{3,2B} \lambda_{2B,2A} \lambda_{2A,1} \lambda_{1,0}}{\left(\left(\lambda_{2A,1}+\lambda_{1,0}\right) \lambda_{2B,2A}+\lambda_{2A,1} \lambda_{1,0}\right) \lambda_{3,2B}+\lambda_{2B,2A} \lambda_{2A,1} \lambda_{1,0}} 
\\
 0 & 
 \Sigma_2 %-\lambda_{2B,2A}-\frac{\lambda_{2B,2A} \lambda_{2A,1}}{\lambda_{2B,2A}+\lambda_{2A,1}}-\frac{\lambda_{2B,2A} \lambda_{2A,1} \lambda_{1,0}}{\left(\lambda_{2A,1}+\lambda_{1,0}\right) \lambda_{2B,2A}+\lambda_{2A,1} \lambda_{1,0}}
 & \lambda_{2B,2A} & \frac{\lambda_{2B,2A} \lambda_{2A,1}}{\lambda_{2B,2A}+\lambda_{2A,1}} & \frac{\lambda_{2B,2A} \lambda_{2A,1} \lambda_{1,0}}{\left(\lambda_{2A,1}+\lambda_{1,0}\right) \lambda_{2B,2A}+\lambda_{2A,1} \lambda_{1,0}} 
\\
 0 & 0 & %-\frac{\lambda_{2A,1} \left(\lambda_{2A,1}+2 \lambda_{1,0}\right)}{\lambda_{2A,1}+\lambda_{1,0}} 
 \Sigma_3 & \lambda_{2A,1} & \frac{\lambda_{2A,1} \lambda_{1,0}}{\lambda_{2A,1}+\lambda_{1,0}} 
\\
 0 & 0 & 0 & %-\lambda_{1,0}
 \Sigma_4 & \lambda_{1,0} 
\\
 0 & 0 & 0 & 0 & 0 
    \end{bmatrix},
\end{align*}
where $\Sigma_i = -\sum_{j=1, i \neq j}^m A^{(s)}_{\RN{2},i,j}$ for $i = 1, \dots, m-1$. 

\myparagraph{Upper-banded (max jump = 2)}
The \textit{upper bidiagonal} parameterization allows for transitions to both the immediate next class and the class two steps ahead in severity. Compared to the generalized Erlang, this formulation introduces additional flexibility, while still maintaining a relatively sparse generator. The generator is given by
\begin{align}
\bm{A}_{\RN{3}}^{(s)} = \begin{bmatrix}
\Sigma_1 & \lambda_{3,2B} & \lambda_{3,2A} & 0 & 0 \\
0 & \Sigma_2 & \lambda_{2B,2A} & \lambda_{2B,1} & 0 \\
0 & 0 & \Sigma_3 & \lambda_{2A,1} & \lambda_{2A,0} \\
0 & 0 & 0 & \Sigma_4 & \lambda_{1,0} \\
0 & 0 & 0 & 0 & 0 \\
\end{bmatrix}, \label{eq:generator_bidiagonal}
\end{align}
where $\Sigma_i = -\sum_{j=1, i \neq j}^m A^{(s)}_{\RN{3},i,j}$ for $i = 1, \dots, m-1$.
The total number of base rate parameters is $2m - 3$. With covariates and warping, the model has $3m - 4 + p$ parameters in total.

\myparagraph{Upper-banded (max jump = 3)}
The fourth parameterization, \textit{upper tridiagonal}, further generalizes the upper bidiagonal by allowing jumps of up to three severity levels ahead. The generator matrix takes the form
\begin{align}
\bm{A}_{\RN{4}}^{(s)} = \begin{bmatrix}
\Sigma_1 & \lambda_{3,2B} & \lambda_{3,2A} & \lambda_{3,1} & 0 \\
0 & \Sigma_2 & \lambda_{2B,2A} & \lambda_{2B,1} & \lambda_{2B,0} \\
0 & 0 & \Sigma_3 & \lambda_{2A,1} & \lambda_{2A,0} \\
0 & 0 & 0 & \Sigma_4 & \lambda_{1,0} \\
0 & 0 & 0 & 0 & 0 \\
\end{bmatrix}, \label{eq:generator_tridiagonal}
\end{align}
where again $\Sigma_i = -\sum_{j=1, i \neq j}^m A^{(s)}_{\RN{4},i,j}$ for $i = 1, \dots, m-1$.
This specification enables modeling large defect escalations, where a defect may deteriorate across multiple classes in a single transition. It introduces more flexibility than the previous structures at the cost of increased cardinality of the parameter space. Specifically, this parameterization includes $3m - 6$ base transition rates, yielding a total of $4m - 7 + p$ parameters with covariates and warping.


\myparagraph{Free upper-triangular}
The last parameterization assumes that a defect can transition to any class of equal or higher severity, without any functional relationship between transition rates. In this setting, the upper triangle of the generator is \textit{free} of structural restrictions. The generator is expressed as
\begin{align}
    \bm{A}_{\RN{5}}^{(s)} =
    \begin{bmatrix} 
    \Sigma_1 & \lambda_{3,2B} & \lambda_{3, 2A} & \lambda_{3, 1} & \lambda_{3, 0} \\
    0 & \Sigma_2 & \lambda_{2B, 2A} & \lambda_{2B, 1} & \lambda_{2B, 0} \\
    0 & 0 & \Sigma_3 & \lambda_{2A, 1} & \lambda_{2A, 0} \\
    0 & 0 & 0 & \Sigma_4 & \lambda_{1, 0} \\
    0 & 0 & 0 & 0 & 0 \\
    \end{bmatrix}, \label{eq:generator_free_upper}
\end{align}
where $\Sigma_i = -\sum_{j=1, i \neq j}^m A_{\RN{5},i,j}^{(s)}$ for $i = 1, \dots, m-1$.
Unlike the reparameterized upper, this formulation does not constrain long-range jumps to have lower rates than neighboring jumps. It introduces more model flexibility at the cost of higher complexity: the number of free parameters increases quadratically with the number of states and linearly with the number of covariates, that is, $(m^2 + m)/2 - 1 +p$ parameters in total with regression and warping.


\FloatBarrier

\subsubsection{Optimum score estimation}\label{sec:optimum_score_estimation}
Let $\bm{\theta}$ denote the set of model parameters to be determined.
For the Markov jump process, $\bm{\theta}$ is the collection of parameters constituted by the base transition rates, $\lambda_{i,j}^{(0)}$, and the coefficients on the covariates, $\bm{\beta}$.
Optimum score estimation is about determining ``good'' model parameters, $\bm{\theta}$, based on maximizing some score function $\mathcal{D}$, i.e.
\begin{align*}
\argmax_{\bm{\theta}}\mathcal{D}_n(\bm{\theta}) \rightarrow \bm{\theta}_0, \quad \text{as } n \rightarrow \infty,   
\end{align*}
\citet{Gneiting2007} argues that the choice of score function, $\mathcal{D}$, should be in accordance with the forecast scoring rule as this improves model performance \citep{friedman1983a,Copas1983} compared to non-tailored estimation practices. 
The overall idea of optimum score estimation is that we choose the measure of the goodness-of-fit in accordance with how we seek to evaluate the performance of the model. 
Given a parametric model, $P_{\bm{\theta}}$, and the recorded data, $Y^{(s)}$, we denote the \textit{metric} (or \textit{reward}) by $D(P_{\bm{\theta}}, Y^{(s)})$, and the mean score by
\begin{align}
    \mathcal{D}_n(\bm{\theta}) = \frac{1}{|\mathcal{S}|} \sum_{s \in \mathcal{S}} D\left(P_{\bm{\theta}}, Y^{(s)}\right). \label{eq:mean_score}
\end{align}
\citet{bisgaard2024A} applies maximum likelihood estimation to determine parameters of the Markov jump process.
The data mean log-likelihood given the recorded jumps of the defect trajectories is 
\begin{equation}
\begin{aligned}
                  \ell\left( \bm{\theta}; \{y_1^{(s)}, y_2^{(s)}, t_1^{(s)},t_2^{(s)}\}_{s \in \mathcal{S}} \right) &= 
                  \frac{1}{|\mathcal{S}|}\sum_{s \in \mathcal{S}} \log \left( P\left(X^{(s)}(t_2^{(s)}) = y_2^{(s)} \mid X^{(s)}(t_1^{(s)}) = y_1^{(s)} \right) \right) \\ 
                  &= \frac{1}{|\mathcal{S}|}\sum_{s \in \mathcal{S}}
                  \log\left(   \left[ \bm{\pi}^{(s)}\left( t_2^{(s)} \right)  \right]_{ y_2^{(s)}} \right)
                  \\
                  &= \frac{1}{|\mathcal{S}|}\sum_{s \in \mathcal{S}}
                  \log\left(   \left[ \bm{e}_{y_{1}^{(s)}}  \exp \left( \bm{A}^{(s)} \tau^{(s)} \right) \right]_{ y_2^{(s)}} \right), \label{eq:likelihood_mjp}
\end{aligned}     
\end{equation}
where $\bm{e}_{y_{1}^{(s)}} \in \mathbb{N}^{1 \times m}$ is $1$ at the state $y_{1}^{(s)}$-index and zero elsewhere (it corresponds to the initial distribution $\bm{\pi}^{(s)}(t_1^{(s)})$). The $y_2^{(s)}$-index on the square brackets indicates that we choose the corresponding element of the transient distribution, $\bm{e}_{y_{1}^{(s)}}  \exp \left( \bm{A}^{(s)} \tau^{(s)} \right)$. 
% state probability of a transition from state $y_1^{(s)}$'th to state $y_{2}^{(s)}$ in the time passed between the two observations, that is $\tau_s = t_2^{(s)} - t_1^{(s)}$. 
Thus, the discrete likelihood is the product of the probabilities of observing the state, $y_2^{(s)}$, after $\tau_s$ time has passed upon observing $y_1^{(s)}$ for all $s \in \mathcal{S}$.
\Cref{eq:likelihood_mjp} is thus an example on a choice of the mean score function in \cref{eq:mean_score}.
Maximum likelihood estimation corresponds to optimum score estimation based on the log-score. Subsequently, MLE is optimal given forecast evaluation wrt. to this exact scoring rule. Provided another scoring rule, MLE might not be suitable for optimal forecast performance. This notion is empirically demonstrated by \citet{Gneiting2005} and \citet{gebetsberger2018a}.
For this study, we apply several probabilistic scoring rules. We describe these in \Cref{sec:error_metrics}. To improve prediction performance wrt. to all metrics, our optimum score estimation is based on maximizing the sum of the mean score functions wrt. their respective scoring rules, i.e.
\begin{align}
     \mathcal{D}_n(\bm{\theta}) = \frac{1}{|\mathcal{S}|} \sum_{s \in \mathcal{S}} \left( D_1\left(P_{\bm{\theta}}, Y^{(s)}\right) + D_2\left(P_{\bm{\theta}}, Y^{(s)}\right) + D_3\left(P_{\bm{\theta}},  Y^{(s)}\right)\right). \label{eq:mean_score_combination}
\end{align}
where $D_i(P_{\bm{\theta}}, Y^{(s)})$ for $i = 1,2,3$ are the three different metrics considered in this study. If deemed relevant, weights can be assigned to each of terms. The following describes the different metrics.

% Model parameters can be determined by maximum likelihood estimation, i.e. maximining the logarithm of the quantity in \cref{eq:likelihood_mjp}, i.e.
% \begin{align*}
%    \widehat{\pmb{\theta}}  \; = \argmax_{\left( \{ \lambda_{i,j}^{(0)} \}_{i,j\in \mathcal{J} | i \neq j}, \bm{\beta} \right)} \sum_{s \in \mathcal{S}} 
%     \log\left(   \left[ \bm{e}_{y_{1}^{(s)}}  \exp \left( \bm{A}^{(s)} \tau^{(s)} \right) \right]_{ y_2^{(s)}} \right). 
% \end{align*}

%T A substantial difference between MLE and optimum RPS estimation is that that optimum RPS estimation evaluates the full distribution whereas MLE relies on evaluating the observed outcome by a logarithmic penalty. Subsequently, it seems straightforward that optimum RPS estimation performs better in terms of the RPS score, as shown by \citet{gebetsberger2018a}; however, it is less clear how the two estimation approaches performs in terms of the Brier score.





\subsection{Metrics} \label{sec:error_metrics}
As $\mathcal{E}$ is a discrete and finite state space, suitable methods for evaluating categorical forecasts must be considered. By \cref{eq:markov_eq_solution}, we model the evolution of the distribution of $\mathcal{E}$, i.e. we produce probabilities associated to each element of the state space.
Opposed to nominal category forecasts, the considered state space of defect severity has a natural ordering, e.g. class $1$ is more severe than class $2A$. Subsequently, it is said to be \textit{ordinal} \citep{broecker2008a,hamill1997,Weijs2010}.
For probabilistic forecasts of an ordinal discrete state space, there are two possibly relevant aspects of verification, i.e. (1) forecast ability in the category corresponding to the outcome of the event, and (2) forecast ability to predict category probabilities \textit{close to} the category corresponding to the outcome of the event \citep{Weijs2010}. 
In the following, we present three different metrics that will be used for both estimation (as described in \Cref{sec:optimum_score_estimation}) and for evaluating model predictions.

%Let $\Tilde{P}^{(s)}_{i}(t)$ denote the predicted probability at time $t$ for outcome $i \in \mathcal{J}$ for defect $s \in \mathcal{S}$. We assume that the $m$ defect classes are mutually exclusive and exhaustive meaning that
% \begin{align*}
%     \sum_{i \in \mathcal{J}} \Tilde{P}^{(s)}_{i}(t) = 1, \quad \forall s \in \mathcal{S}, \forall t \in \mathbb{R}_+.
%\end{align*}
%We assume availability of discrete data, i.e. observations of trajectorial endpoints $Y^{(s)} = \{ y_1^{(s)} = X^{(s)}(t_1^{(s)}), \; y_2^{(s)} = X^{(s)}(t_2^{(s)}) \}_{\forall s \in \mathcal{S}}$ with $y_1^{(s)}$ and $y_2^{(s)}$ being the observed states and $t_1^{(s)}$ and $t_2^{(s)}$ their respective time stamps, and $X^{(s)}$ the Markov process. We thus consider the endpoints of the defect trajectories as ''outcomes`` or ''targets`` for prediction evaluation purposes, that is the finite set of discrete values, $\{y_2^{(s)} \}_{\forall s \in \mathcal{S}}$.


\subsubsection{Log Score}
The log-score is a widely used scoring rule; optimum score estimation based on the log-score corresponds to maximum likelihood estimation. The log-score evaluates the logarithm of the probability assigned to the observed outcome. In the case of the current study, the mean log-score can be expressed as
\begin{align}
    \overline{\text{LogS}} = \frac{1}{\vert \mathcal{S} \vert} \sum_{s \in \mathcal{S}} \log \left( \left[ \bm{\pi}^{(s)}\left( t_2^{(s)} \right)  \right]_{ y_2^{(s)}} \right).  \label{eq:log_Score}
\end{align}
\Cref{eq:log_Score} is thus equivalent to \cref{eq:likelihood_mjp}.
It rewards forecasts solely based on the probability assigned to the observed category, ignoring the probabilities assigned to incorrect ones. Therefore, it is \textit{insensitive to distance} in the state space.  


\subsubsection{Brier score} 
\citet{Brier1950} proposed a scoring rule for multicategorical predictive distributions, which in our case is given as the mean of the sum of the squared distances of the computed transient distributions and the observed endpoints of the defect trajectories, i.e.
\begin{align}
    \overline{\text{BS}} = \frac{1}{\vert \mathcal{S} \vert} \sum_{s \in \mathcal{S}} \sum_{i \in \mathcal{J}} \left( \bm{\pi}^{(s)}_i\left( t_2^{(s)} \right) -\mathds{1}_{ \{i \; = \; y_2^{(s)} \} } \right)^2, \label{eq:Brier_Score}
\end{align}
where $\mathds{1}_{ q }$ is an indicator which is one if $q$ is true and zero otherwise.
Like the log-score, the Brier score treats classes not corresponding to the observed outcome equally, that is by summing over the quadratic errors, making it also insensitive to distance. A perfect forecast is associated with a zero contribution to the Brier score for each defect, $s$. Therefore, it is a loss function.


\subsubsection{Ranked Probability Score}
The mean ranked probability score (RPS) is given by
\begin{align}
    \overline{\text{RPS}} =\frac{1}{\vert \mathcal{S} \vert} \sum_{s \in \mathcal{S}} \sum_{k = 1}^m\left( \sum_{i = 1}^k\bm{\pi}^{(s)}_i\left( t_2^{(s)} \right)  - \sum_{i = 1}^k\mathds{1}_{ \{ i \; = \; y_2^{(s)} \} } \right)^2 \label{eq:RPS}
\end{align}
as the mean of squared distances between the cumulative transient distribution and the cumulative observed outcomes \citep{Epstein1969}.
Compared to the log-score and the Brier score, RPS punishes probability masses in false outcomes milder if the wrong predictions are located \textit{closer} to the observed outcome. Thus, RPS introduces an ordering to inaccurate predictions in the sense that some \textit{wrong} predictions can be less wrong.
For the purpose of degradation dynamics where there is a structural ordering wrt. defect severity, it can be argued that RPS is more suitable than the log-score or the Brier Score, if forecasts are interpreted as beliefs about track safety levels, which correlate with crack size \citep{UIC_2002_defecthandbook}. 
According to \citet{Czado2009}, an advantage with RPS is that it enables direct comparison between predictive distributions and point predictions, as it has shown to generalize the mean absolute error \citep{Hersbach2000}.

\subsubsection{Metrics recap} \label{sec:metrics_recap}
With these three scoring rules, we have outlined a collection of metrics to be used for both model estimation (via optimum score estimation) and the evaluation of model predictions. The choice of metric depends on the problem at hand and which forecast properties are preferred by the end-user. The present study outlines a general model framework that can be adapted to handle any of the three metrics or a combination of them.
With that said, the log-score and the Brier score are inherently suited for verification of forecast ability in the category corresponding to the outcome of the event (nominal state space), while RPS is better suited for evaluating forecasts for ordinal state spaces. Therefore, we suggest that more emphasis should be placed on the RPS for studies related to degradation. However, there might be reasons to differ from that general suggestion, for instance if safety precautions differ substantially from one category to the other.
When used for estimation, loss functions like the Brier score and RPS are minimized (hence enter with a negative sign in the score function $\mathcal{D}$), while the log-score, being a utility function, is maximized.


\subsection{Computational remarks} \label{sec:computation_remarks}
We have implemented the Markov jump process in \texttt{C++} integrated with \texttt{R} via \texttt{Rcpp} \citep{Rcpp}. The repository can be accessed on \url{https://github.com/dobbeltgaard/MarkovJump.git}. The implementation allows the user to specify which metric(s) to use for optimum score estimation (log score, Brier score, ranked probability score). Additionally, the implementation allows the user to specify one of the three parameterizations of the generator from \Cref{sec:generator_params}. 
Prediction is performed by applying the solution to the Kolmogorov differential equation in \cref{eq:markov_eq_solution}. All aforementioned scoring rules are implemented for prediction evaluation. 
The model implementation accommodates three calculation methods for the transient distribution of the Markov chain. The first is a Padé approximation of the matrix exponential. The second method is to calculate the transient distribution via \textit{uniformization} (or \textit{randomization}). Uniformization provides an alternative representation of a Markov jump process such that all holding times distributions have the same intensity parameter \citep{bladt_nielsen}. We refer to \citet{sherlock2021a} for a detailed algorithmic description. The third method is by applying the eigendecomposition
\begin{align*}
    \exp{\left( \bm{A}u \right) } = \bm{U}  \exp{\left( \bm{\Lambda}u \right) } \bm{U}^{-1}, 
\end{align*}
which reduces the numerical effort substantially since the matrix exponential to a diagonal matrix is the same as taking the exponential to each diagonal element.
We refer to \citet{bisgaard2024A} for algebraic expressions for the eigenspaces $\bm{U}$ and $\bm{U}^{-1}$, in case of the generalized Erlang parameterization. In case of the other two generator parameterizations, we apply numerical schemes for the eigendecompositions using the \texttt{Eigen} library in \texttt{C++} \citep{eigenweb}.

We have tested the speed of the score computation with respect to the different score options and the different methods for calculating the transient distribution of the Markov chain, using the \texttt{bench} library \citep{R_bench}. We compare our \texttt{C++} implementation with a standard (``slow'') \texttt{R} implementation based on a Padé approximation of the transient distribution. \Cref{tab:mjp_score_computation_speed} displays the results of the speed investigation.
\begin{table}[htb!]
    \centering
    \caption{Relative speed of Markov jump score computation wrt. calculation method of transient distribution and score metric. Uniformization is performed with maximum truncation error on $2^{-52}$ (see \citet{sherlock2021a} for details). The base reference for the speed investigation is the ''\texttt{C++} w. eigendecomposition`` which has a median run time on $1.85$ ms.}
    \label{tab:mjp_score_computation_speed}
\begin{tabular}{lccc}
\toprule
  & Log score & Ranked probability score & Brier score \\
\midrule
\texttt{R} w. Padé & 154.03 & 98.71 & 150.15\\
\texttt{C++} w. Padé & 3.26 & 3.40 & 3.42\\
\texttt{C++} w. uniformization & 6.77 & 6.84 & 6.78\\
\texttt{C++} w. eigendecomposition & \textbf{1.00} & 1.10 & 1.07\\
\bottomrule
\end{tabular}
\end{table}
As seen, there is a substantial speedup going from a standard \texttt{R} implementation to a corresponding \texttt{C++} implementation which is approximately 50 times faster. The calculation method for the transient distribution does impact the computation speed. If the generator matrix is diagonalizable, the eigendecomposition provides a highly efficient (and accurate) way to calculate the transient distribution. Our implementation checks if this is the case, and if not, uniformization is applied as outlined by \citet{sherlock2021a}. Although uniformization can be slower than the Padé approximation, it has the advantage that the user controls the truncation error. For the \texttt{C++} implementations, there are only minor differences in computation speed wrt. which score is applied. This is expected as there is only a small difference in the number of operations required to compute the various scores. However, we observe a substantial difference of the various \texttt{R} implementations. The difference between them is most likely explained by the fact that the incorporated \texttt{cumsum} function in the RPS implementation can benefit from CPU cache better since it processes data in a sequential manner, making it more cache-friendly. Additionally, the \texttt{cumsum} operation has some inherent optimization that might make it faster than a simple element-wise difference and squaring.



\subsection{Reference predictors} \label{sec:reference_forecasts}
The use of reference predictors (or alternatively \textit{baseline predictors}) is essential in forecast studies, as it allows comparison between various methods. Benchmarking a proposed model with reference predictors reveals the gains of introducing it wrt. some chosen score/metric (see \Cref{sec:error_metrics}). For instance, it enables quantification of what is gained by introducing complex model structures with e.g. temporal and state dependencies or covariates.
In the following, we describe the reference predictors that are used to benchmark the Markov jump process from \Cref{sec:mjp}.
We note that we have used the ''trick`` of introducing simple state dependence to these reference predictors so that they only assign probability mass to states that are possible for $X^{(s)}(t)$ to reach given its current state. For instance, if the most recent observation of $X^{(s)}(t)$ is a class $2A$ defect, it does not make sense to predict any probability mass in the classes $3$ or $2B$, as this study considers defect trajectories subject to non-reversible degradation.  
We do this to ensure that the reference predictors used to benchmark the Markov jump process are competitive and that the Markov jump process is compared to predictors that constitute actual reference predictability.


\subsubsection{Uniform predictor}
Predicting the uniform distribution is a type of \textit{preferenceless} naïve forecast method, i.e. it predicts equal probability masses over the entire state space. Subsequently, each state is assigned a probability of $1/m$. As defect propagation is state dependent and assumed irreversible, we correct the uniform predictor so that the entire probability mass is divided between the states that are possible to reach given the starting point $y_1^{(s)}$ for each defect. 
Thus, the corrected uniform predictor is given by 
\begin{align}
    \widehat{\bm{\pi}}^{(s),\text{Uni.}}_i \left( y_1^{(s)} \right) = \frac{\mathds{1}_{ \{  i \; \geq \; y_1^{(s)}  \} }   }{\sum_{j \in \mathcal{J}} \mathds{1}_{ \{ j \; \geq \; y_1^{(s)} \} }   }, \quad \forall s \in \mathcal S. \label{eq:uniform_predictor}
\end{align}
The uniform predictor in \cref{eq:uniform_predictor} can be regarded as a variation of a random classifier, i.e. it is preferenceless in the reachable states.



\subsubsection{Empirical distribution predictor}
Predicting the empirical distribution is another naïve approach to forecasting. It is more competitive given data with many repeated events. For this study, we compute the empirical distribution of the initial state of the defect trajectories, i.e. $y_1^{(s)}$. Given the distribution of the recorded $y_1^{(s)}$ for $s \in \mathcal{S}$, the empirical predictor assigns corresponding probability masses to the state space for the future, i.e.
\begin{align}
    \widehat{\bm{\pi}}^{(s),\text{Empi.}}_i \left( y_1^{(s)},y_2^{(s)} \right) = \frac{\sum_{s \in \mathcal{S}} \mathds{1}_{ \{  i \; = \; y_1^{(s)} \; \wedge \; j \; = \; y_2^{(s)}   \} }   }{\sum_{s \in \mathcal{S}} \mathds{1}_{ \{  i \; = \; y_1^{(s)}    \} }   }, \quad \forall j \in \mathcal{J}.  \label{eq:empirical_distr}
\end{align}
% \begin{align}
%     \widehat{\bm{\pi}}^{(s),\text{Empi.}}_i \left( y_1^{(s)} \right) = \begin{cases}
%         0, \quad i \; < \; y_1^{(s)} \\
%         \frac{1}{|\mathcal{S}|}\left( \sum_{s \in \mathcal{S}} \left( \mathds{1}_{ \{ i \; = \; y_1^{(s)} \} } +
%         \frac{\mathds{1}_{ \{ i \; < \; y_1^{(s)} \} }}{m - y_1^{(s)} + 1}  \right) \right)
%         , \quad \text{else}.
%     \end{cases}
%     \label{eq:empirical_distr}
% \end{align}
%As described earlier, we add $\frac{\mathds{1}_{ \{ i \; < \; y_1^{(s)} \} }}{m - y_1^{(s)} + 1} $ to the reachable states to correct for probability masses assigned to states that are unreachable for $X^{(s)}$ given $y_1^{(s)}$.



\subsubsection{Cumulative link model} \label{sec:cumulative_link_model}
The cumulative link model is a static regression model that predicts probabilities for an ordinal state space. Opposite to the case of the uniform predictor and the empirical distribution predictor, it cannot be said to be a naïve method. It is included to the collection of reference predictors to compare the benefits of the dynamic model setting of the Markov jump process.
The cumulative probabilities are defined as
\begin{align*}
     P\left( X^{(s)} \leq j \; \vert \; \bm{z}'^{(s)} \right) = 
     \pi_1^{(s)}(\bm{z}'^{(s)}) + \dots + \pi_j^{(s)}(\bm{z}'^{(s)}), \quad j = 1, \dots, m, 
\end{align*}
where $\pi^{(s)}_j\left(\bm{z}'^{(s)}\right) = P \left(  X^{(s)} = j \; \vert \; \bm{z}'^{(s)} \right)$ with $\sum_{j = 1}^m \pi^{(s)}_j(\bm{z}'^{(s)}) = 1$ for $s \in \mathcal{S}$. $\bm{z}'^{(s)}$ is the set of covariates (regression variables) for each defect $s \in \mathcal{S}$. Note the staticity is indicated by writing $X^{(s)}$, while we in \Cref{sec:mjp} denoted the process by $X^{(s)}(t)$ with explicit time dependence.
$G^{-1}$ denotes a link function, which is the inverse of some continuous cumulative density function, $G$.
The cumulative link model is then given by
\begin{align}
    G^{-1} \left(  P\left( X^{(s)} \leq j \; \vert \; \bm{z}'^{(s)} \right) \right)
    = \alpha_j + \bm{\beta}^\top \bm{z}'^{(s)}, \quad j = 1, \dots, m-1, \; s \in \mathcal{S}, \label{eq:cum_link_model}
\end{align}
where $\alpha_j$ is an intercept for the $j$ cumulative link, and $\bm{\beta}^\top$ are the coefficients for the covariates, $\bm{z}'^{(s)}$. 

As seen from \cref{eq:cum_link_model}, the cumulative probability of an event, $j$, is obtained by the linear predictor, $\alpha_j + \bm{\beta}^\top \bm{z}'^{(s)}$, through some link, $G$.
The regression setting in \cref{eq:cum_link_model} is somewhat related to that of the regression in \cref{eq:trans_rates_covariates} as there is some base rate (or intercept) that is scaled (or added) by the covariate contribution, all mapped through some transformation of exponential family (link function). However, the covariate set, $\bm{z}^{(s)}$, in \cref{eq:trans_rates_covariates}, is only a subset of the covariates used in \cref{eq:cum_link_model}. As the Markov jump process inherently captures temporal dynamics and state dependency via \cref{eq:markov_property}, there is no need to include the time between observations, $\tau^{(s)}$, and the initial defect observation $y_1^{(s)}$ in the covariates $\bm{z}^{(s)}$. 
The cumulative link model is static in terms of temporal dependency. Further, it has no intrinsic formulation of state dependence. This implies that some of the recorded endpoint data of the defect trajectories, $Y^{(s)}$, need to be incorporated through the covariates, i.e. the time between observations, $\tau^{(s)}$, and the initial defect observation $y_1^{(s)}$. Subsequently, $\bm{z}'^{(s)}$ contains this information in addition to the rail characteristics described in \Cref{sec:data}. Therefore, it is denoted by a prime to indicate the difference between the covariates in \cref{eq:trans_rates_covariates}.


The link function specifies the distribution of an unobserved continuous variable, of which only crude ordinal measurements, $Y^{(s)}$, are recorded. If $G(\cdot)$ is the standard logistic distribution, then $G^{-1}(\cdot)$ is the logit function, leading to a \textit{proportional odds logistic regression model} \citep{mccullagh1980a}, where the log-odds of being in category $j$ or below differ by a constant across categories.
If $G(\cdot)$ is the standard normal distribution, then $G^{-1}(\cdot)$ is the probit function, and the cumulative link model generalizes the binary probit model \citep{agresti2012}.
If $G^{-1}(\cdot)$ is the complementary log-log function, the model corresponds to a proportional hazards model for grouped survival times \citep{Christensen2023}.
In the prediction study in \Cref{sec:results}, we test the logit, probit, and complementary log-log link functions to evaluate which yields the highest prediction accuracy.


By the law of total probability for a discrete random variable, we have that 
\begin{align}
     \pi^{(s)}_j\left(\bm{z}'^{(s)}\right) = 
     P\left( X^{(s)} = j \; \vert \; \bm{z}'^{(s)} \right) = 
    \begin{cases}
    P\left( X^{(s)} \leq j \; \vert \; \bm{z}'^{(s)} \right), \quad j = 1 \\
     P\left( X^{(s)} \leq j \; \vert \; \bm{z}'^{(s)} \right) - P\left( X^{(s)} \leq j -1 \; \vert \; \bm{z}'^{(s)} \right), \quad \text{else}.   
    \end{cases} \label{eq:cum_link_probs}
\end{align}
The likelihood function of the recorded jumps is then obtained from the probability mass function of the categorical distribution (the single trial multinomial distribution corresponding to the multivariate extension of the Bernoulli distribution), i.e.
\begin{equation*}
\begin{aligned}
     L\left( \bm{\theta}; y_2^{(1)}, \dots, y_2^{(|\mathcal{S}|)} \right) = \prod_{s \in \mathcal{S}} \prod_{j = 1}^m \left( \pi^{(s)}_j\left(\bm{z}'^{(s)}\right) \right)^{\mathds{1}_{ \{ j \; = \; y_2^{(s)}  \} }}, 
\end{aligned} \label{eq:likelihood_cumlink}
\end{equation*}
where $\mathds{1}_{ \{ q \} }$ is an indicator which is 1 if the event $q$ occurs and zero otherwise, and $\bm{\theta}$ denotes the set of model parameters (that is $\{\alpha_j\}$ and $\bm{\beta}$). The model parameters are determined by maximum likelihood estimation taking the logarithm to the likelihood function, i.e.
\begin{align*}
   \widehat{\pmb{\theta}}  \; = \argmax_{\left( \{ \alpha_{j} \}_{j = 1}^{m-1}, \bm{\beta} \right)} \sum_{s \in \mathcal{S}} \sum_{j = 1}^m 
   \mathds{1}_{ \{ j \; = \; y_2^{(s)} \} } \log\left( \pi^{(s)}_j\left(\bm{z}'^{(s)}\right) \right), 
\end{align*}
where the $\pi^{(s)}_j\left(\bm{z}'^{(s)}\right)$ are defined in \cref{eq:cum_link_probs} related to the cumulative link functional in \cref{eq:cum_link_model}.
The model is estimated with the \texttt{polr} function from the \texttt{MASS} library \citep{mass_R}.
With determined model coefficients, the cumulative link model can be applied for prediction. By model definition, probabilities are forecasted. Given a set covariates $\bm{z}'^{(s)}$ for a defect $s \in \mathcal{S}$, that is the rail characteristics described in \Cref{sec:data}, and the time between observations, $\tau^{(s)}$, and the initial defect observation $y_1^{(s)}$, we compute the probabilities for each category by \cref{eq:cum_link_probs} via the cumulative links in \cref{eq:cum_link_model}.

